\documentclass[11pt,a4paper]{article}
\usepackage[T1]{fontenc}
\usepackage[utf8]{inputenc}
\usepackage[margin=27mm]{geometry}
\usepackage{amsmath,amssymb,booktabs,graphicx,microtype}
\usepackage[numbers,sort&compress]{natbib}
\usepackage[hidelinks]{hyperref}

% A starting layout, not a completed paper. Replace with the supplied venue class when required.
\title{Research question}
\author{}
\date{}

\begin{document}
\maketitle
\begin{abstract}
State the question, method, observed result and limitation. Do not imply an
experiment was conducted when this document is a proposal or literature study.
\end{abstract}

\section{Introduction}
Define the problem, scope and contribution using verified evidence.

\section{Related work}
Compare the primary sources actually read; add their verified bibliography entries.

\section{Method}
Describe the data, assumptions, procedure and evaluation so another person can
reproduce the work. State what was executed and what is only proposed.

\section{Results}
Report measured results with units and uncertainty. Keep illustrative data
explicitly labeled. Do not fabricate results to complete this template.

\section{Discussion and limitations}
Explain what follows from the evidence and what remains uncertain.

\section{Conclusion}
Answer the research question to the extent the evidence permits.

% Enable the bibliography after adding verified sources and citations.
% \bibliographystyle{plainnat}
% \bibliography{references}
\end{document}
